\subsection{The range of the Pfaffian bound}

\begin{proposition}\label{prop:pfaffian-range}
Let \(A/\mathbb F_p\) have dimension \(2m\) and a prime-to-\(p\)
polarization, with \(p>3\). Suppose an imaginary quadratic field \(E\)
acts geometrically, its action is integral and maximal at \(p\), and
\(p\) is unramified in \(E\). Suppose Rosati and Frobenius both act on
\(E\) by conjugation and the two tangent multiplicities are \((m,m)\).
Write \(\pi\) for Frobenius and \(D=\det_E(1+\pi^2/p)\). If \(A\)
is nonordinary, then \(v_p(D)\ge 2-m\), with \(v_p(0)=+\infty\).
For every \(m\ge3\) and every prime \(p>3\) with \(p\equiv2\pmod3\),
there is a principally polarized example with \(E=\Q(i)\) and
\(D=p^{2-m}\).
\end{proposition}

\begin{proof}
The two blocks of \(\pi^2\) are conjugate by \(\pi\), so the two
embeddings of \(D\) agree. Over an unramified Witt ring splitting the
coefficient action, let \(B\) be the off-diagonal block of \(\pi\) in
dual integral bases. The Frobenius lattice maps have \(m\) elementary
divisors \(1\) and \(m\) elementary divisors \(p\), by the two tangent
multiplicities. Hence \(v_p(\det B)=m\). The same matrix identity as in
the fourfold gives
\[
 D=\frac{\operatorname{Pf}(B^{\mathsf T}-B)^2}{\det B}.
\]
In each block the eigenvalue valuations of \(\pi^2/p\) are reciprocal
pairs \(\pm a_j\) with \(0\le a_j\le1\), for \(1\le j\le m\).
If the reduction is nonordinary, \(\sum_j a_j<m\). Thus
\(v_p(D)\ge-\sum_j a_j>-m\), and the Pfaffian gives
\(v_p(D)\equiv m\pmod2\) when \(D\ne0\). This yields the bound.

For sharpness, take an ordinary elliptic curve \(O/\mathbb F_p\) with
Frobenius polynomial \(X^2-X+p\) and a supersingular elliptic curve
\(T/\mathbb F_{p^2}\) with Frobenius \(p\zeta_3\). Both exist by
Honda--Tate \cite[Theorem~1]{tate}. The first Weil number has a quadratic
imaginary field and local slopes zero and one. For the second, the
field is \(\Q(\zeta_3)\), in which \(p\) is inert. The sole local
invariant at \(p\) is \(2\cdot(1/2)=0\bmod1\), so its Schur index is
one and its isogeny class has dimension one. Put
\[
 S=\operatorname{Res}_{\mathbb F_{p^2}/\mathbb F_p} T,
 \qquad U=O^{m-2}\times S.
\]
The polynomial of \(S\) is \(X^4+pX^2+p^2\), by the two-step
Frobenius action on the geometric factors. Their product polarization
descends and is principal. Let \(U'\) be the quadratic twist of \(U\)
by \([-1]\) and put \(A=U\times U'\), with the product polarization.
Under the geometric identification with \(U^2\), let \(i\) act by
\((x,y)\mapsto(-y,x)\). Its Rosati adjoint is \(-i\), and its tangent
multiplicities are \((m,m)\). Frobenius acts as
\(\operatorname{diag}(\pi_U,-\pi_U)\), so conjugates \(i\).

On either coefficient block the square of Frobenius acts as
\(\pi_U^2\). If \(\alpha,\beta\) are the roots for \(O\), its
contribution to \(D\) is
\((1+\alpha^2/p)(1+\beta^2/p)=(\alpha+\beta)^2/p=1/p\).
For \(S\) the contribution is
\((1+\zeta_3)^2(1+\zeta_3^2)^2=1\). Therefore \(D=p^{2-m}\),
while the supersingular factor makes \(A\) nonordinary.
\end{proof}

In rank four the bound makes the norm integral at nonordinary primes.
In higher even rank, the examples have nonintegral norm with the same
local lattice and Weil constraints. They establish sharpness for the
finite-field bound.
